\section{本章小结}

本章介绍了TIDE-Swarm的三项核心机制。
TIDE-Anchor通过Context Ledger、分层记忆与原文归档维护任务信息，
结合预算检查、漂移检测和检查点支持持续执行；
TIDE-Mesh根据任务需求选择异构成员，
通过转交、替补与结构化通信组织协作；
TIDE-Placement结合任务画像、可信域策略和资源状态选择推理服务，
并利用执行反馈进行有限重调度。

神经符号协同通过计划中间表示、显式校验、反例修订和质量反馈贯穿上述机制，将任务语义、动态拓扑与资源决策接入同一可观测闭环。其设计优势在于稳定保存目标与证据、按需组织异构能力、保护约束变化并利用失败反馈调整后续执行；核心算法、适用范围及计算开销共同构成平台实现和实验评估的依据。
下一章介绍相关机制在平台中的集成及典型任务应用。


